\documentclass[12pt]{article}
\usepackage[utf8]{inputenc}
\usepackage[brazil]{babel}
\usepackage{parskip}
\usepackage{indentfirst}
\usepackage{titlesec}
\usepackage{geometry}
\usepackage{setspace}
\usepackage[colorlinks=true, urlcolor=blue]{hyperref}

\geometry{a4paper, margin=2.5cm}

\titleformat{\section}
  {\bfseries\fontsize{13}{15}\selectfont}
  {\thesection}{1em}{}
\titlespacing{\section}{0pt}{12pt}{6pt}

\titleformat{\subsection}
  {\bfseries\fontsize{12}{14}\selectfont}
  {\thesubsection}{1em}{}
\titlespacing{\subsection}{0pt}{12pt}{4pt}

\setlength{\parindent}{1.27cm}

\begin{document}
\thispagestyle{empty}

\begin{center}
  \vspace{12pt}
  {\fontsize{16}{19}\bfseries Agente de IA para revisão de código Python com validação PEP 8 automatizada\par}
\end{center}

\vspace{12pt}
\begin{center}
  {\fontsize{12}{14}\bfseries
    Felipe U. G. de Sousa; Gustavo N. Siqueira; Thomaz de S. Scopel; Vinicius S. Cappatti\par}
\end{center}

\vspace{12pt}
\begin{center}
  {\fontsize{12}{14}\selectfont
    Universidade Presbiteriana Mackenzie\\
    Rua da Consolação, 930 --- São Paulo, SP, Brasil\par}
\end{center}

\vspace{6pt}
\begin{center}
  {\fontfamily{pcr}\fontsize{10}{12}\selectfont
    10418415@mackenzista.com.br, 10419057@mackenzista.com.br, 10417183@mackenzista.com.br, 10418266@mackenzista.com.br\par}
\end{center}

\vspace{6pt}
\begin{list}{}{
  \setlength{\leftmargin}{0.8cm}
  \setlength{\rightmargin}{0.8cm}
}
\item
  {\fontsize{12}{14}\selectfont
    \textbf{Resumo.} Este trabalho apresenta um agente de inteligência artificial para revisão de código Python orientada a boas práticas. A solução combina análise estática determinística com a biblioteca \texttt{pycodestyle} e geração de feedback textual por meio de um modelo de linguagem executado localmente via \texttt{ollama}. O pipeline foi implementado para analisar arquivos individuais, diretórios locais e repositórios públicos do GitHub. Em uma base de 25 arquivos de teste, o componente de validação PEP 8 identificou 217 violações em 12 arquivos e nenhum erro em 13 arquivos. A resposta gerada pelo modelo foi condicionada pelas violações coletadas para reduzir alucinações e manter rastreabilidade entre entrada e saída. Como contribuição, o projeto entrega uma arquitetura reprodutível, com foco em explicação didática de erros de estilo, privacidade dos dados e apoio ao aprendizado de programação.\par}
\end{list}

\begin{list}{}{
  \setlength{\leftmargin}{0.8cm}
  \setlength{\rightmargin}{0.8cm}
}
\item
  {\fontsize{12}{14}\selectfont
    \textbf{Abstract.} This work presents an AI agent for Python code review guided by good programming practices. The solution combines deterministic static analysis through \texttt{pycodestyle} with natural-language feedback generation using a local language model via \texttt{ollama}. The pipeline supports individual files, local directories, and public GitHub repositories. In a 25-file benchmark, the PEP 8 validation component detected 217 violations in 12 files, while 13 files contained no violations. The model response is grounded on the collected violations to reduce hallucinations and improve traceability. The main contribution is a reproducible architecture focused on didactic feedback, data privacy, and programming education support.\par}
\end{list}

\section{Introdução}

A revisão de código é uma prática central no desenvolvimento de software, porém continua sendo uma atividade custosa e, em ambientes educacionais, nem sempre disponível com a frequência necessária para orientar estudantes iniciantes. Em disciplinas introdutórias de programação, é comum que problemas de estilo e legibilidade persistam durante boa parte do ciclo de aprendizagem, mesmo quando os algoritmos estão funcionalmente corretos. Esse cenário motiva a adoção de ferramentas que ofereçam feedback imediato e padronizado sobre a qualidade do código.

No ecossistema Python, a PEP 8 estabelece convenções amplamente aceitas para escrita de código legível e consistente. Apesar de sua relevância, a aplicação prática dessas regras depende de ferramentas de análise estática ou da experiência de revisores humanos. Assim, a proposta deste projeto é construir um agente de IA capaz de revisar código Python com base em resultados objetivos produzidos por \texttt{pycodestyle}, transformando erros técnicos em orientações didáticas compreensíveis.

A justificativa do trabalho está na integração entre dois tipos complementares de inteligência computacional: a validação determinística, que reduz ambiguidades na identificação de violações de estilo, e o modelo de linguagem, que amplia a utilidade pedagógica do resultado ao explicar causas e correções. O objetivo principal é disponibilizar uma ferramenta de linha de comando que receba código Python de diferentes fontes, execute validação PEP 8 automaticamente e gere um relatório estruturado de revisão. Como objetivo específico, busca-se preservar a privacidade do código analisado por meio de execução local do modelo.

Este artigo está organizado da seguinte forma: a Seção 2 apresenta a fundamentação teórica e os trabalhos relacionados; a Seção 3 descreve metodologia, dataset e resultados esperados; a Seção 4 discute resultados obtidos e aspectos éticos; a Seção 5 sintetiza as conclusões; e a Seção 6 lista as referências bibliográficas.

\section{Fundamentação Teórica}

\subsection{PEP 8, análise estática e qualidade de código}

A PEP 8 define recomendações para estilo de código Python com foco em legibilidade, consistência e manutenção de longo prazo. Regras como espaçamento entre operadores, uso adequado de linhas em branco e limite de tamanho de linha reduzem ruído cognitivo e tornam revisões mais eficientes. Ferramentas como \texttt{pycodestyle} automatizam a verificação dessas regras por meio de inspeção lexical e sintática leve, retornando códigos de violação padronizados (por exemplo, E225, E231, E302).

Do ponto de vista de engenharia de software, a análise estática tem a vantagem de produzir resultados reprodutíveis e auditáveis. Em ambientes acadêmicos, essa propriedade é especialmente importante para avaliação formativa, pois permite comparar desempenhos ao longo do tempo com critérios estáveis.

\subsection{Agentes com LLM para revisão de código}

Modelos de linguagem de grande porte (LLMs) ampliaram o escopo da revisão automática de código ao introduzir capacidade de explicação em linguagem natural e sugestões contextualizadas. Entretanto, sem ancoragem em sinais objetivos, esses modelos podem apresentar inconsistências, como contagens incorretas ou recomendações não relacionadas ao problema detectado. Uma estratégia recorrente na literatura recente é combinar componentes simbólicos e determinísticos com geração textual, delegando a identificação do erro a um verificador e a explicação ao modelo generativo.

Esse acoplamento é adotado neste projeto: o \texttt{pycodestyle} atua como fonte de verdade para violações PEP 8 e a LLM atua como camada de interpretação pedagógica. A combinação reduz alucinações e melhora a confiabilidade do feedback final.

\subsection{Trabalhos relacionados}

Simões e Venson (2024) avaliaram LLMs na classificação de qualidade de código Java em comparação com SonarQube e observaram divergências entre abordagens automáticas baseadas em métricas e respostas gerativas. O estudo reforça a necessidade de critérios explícitos quando se utiliza IA para avaliação técnica.

Wang et al. (2024) discutem arquiteturas de agentes baseados em LLM com componentes de memória, planejamento e ação orientada por ferramentas externas, abordagem alinhada ao desenho deste trabalho, no qual a LLM consulta indiretamente saídas de \texttt{pycodestyle} para embasar respostas.

Plaat et al. (2025) apresentam uma revisão sobre modelos agenticos e argumentam que sistemas híbridos, com instrumentos externos de verificação, tendem a produzir respostas mais robustas em tarefas estruturadas. No contexto deste projeto, isso sustenta o uso da validação PEP 8 como etapa obrigatória antes da geração textual.

\section{Metodologia e Resultados Esperados}

\subsection{Arquitetura e pipeline}

O sistema foi implementado em Python com execução por linha de comando e três modos de entrada: arquivo local (\texttt{-f}), diretório local recursivo (\texttt{-d}) e repositório público do GitHub (\texttt{-g}). A execução inicia com o parser de argumentos, que impõe exclusividade entre as fontes de entrada e captura parâmetros de controle como modo verboso e limites de coleta no GitHub. Após o recebimento dos argumentos, o programa resolve a origem dos dados: no modo arquivo, lê os caminhos informados individualmente; no modo diretório, percorre a árvore de arquivos de forma recursiva; no modo GitHub, consulta a API para listar e baixar conteúdos de código dentro dos limites configurados de quantidade e tamanho.

Concluída a aquisição, os conteúdos são normalizados em um dicionário \texttt{\{nome\_arquivo: conteudo\}}, permitindo tratamento uniforme das três modalidades de entrada. Em seguida, ocorre a filtragem estrita por extensão \texttt{.py} e a exclusão de diretórios não relevantes para análise estática (por exemplo, \texttt{.venv}, \texttt{.git} e \texttt{\_\_pycache\_\_}). Essa etapa evita ruído no processo e garante que apenas artefatos Python válidos avancem para validação.

Na sequência, cada conteúdo é submetido ao \texttt{pycodestyle} por meio de arquivo temporário, estratégia necessária porque a biblioteca opera naturalmente sobre caminhos de arquivo. O validador executa as verificações PEP 8 e retorna uma coleção estruturada de violações por arquivo, contendo linha, coluna, código e mensagem. As ocorrências são ordenadas para estabilizar a saída e facilitar reprodutibilidade entre execuções. Ainda nessa fase, o programa calcula métricas agregadas, como total de violações e distribuição por categoria, utilizadas tanto na saída de terminal quanto no contexto enviado ao modelo.

Com os resultados determinísticos consolidados, o sistema monta o contexto de inferência por arquivo combinando código-fonte integral, lista de violações e resumo por categoria de erro. Esse material é interpolado em um prompt com restrições explícitas de confiabilidade, instruindo o modelo a não inventar contagens ou códigos não presentes na entrada e a sinalizar recomendações fora do escopo do \texttt{pycodestyle} como sugestões extras. O prompt final é então enviado ao modelo local via \texttt{ollama}, preservando confidencialidade do código ao evitar transmissão para serviços externos.

Por fim, a resposta gerada é persistida em \texttt{resposta\_llama.md}, acompanhada da identificação da fonte analisada. O fluxo também imprime no terminal um resumo operacional com quantidade de arquivos, número de violações por arquivo e estatísticas globais de execução, permitindo auditoria rápida do processamento. Como critério de confiabilidade, a ferramenta foi configurada para que o relatório textual considere exclusivamente as violações detectadas na validação PEP 8, marcando qualquer recomendação fora desse escopo como complementar.

\subsection{Dataset e análise exploratória}

O conjunto de avaliação contém 25 arquivos Python de perfil didático (funções, classes e exercícios básicos). Desses, 12 arquivos possuem violações PEP 8 intencionais e 13 não possuem violações. O gabarito foi preparado previamente e disponibilizado em \texttt{tests/gabarito/gabarito.md}.

Na análise exploratória inicial, observou-se predominância de erros de espaçamento e separação visual entre blocos de código, com maior incidência nas categorias E231, E225 e E302. Essa distribuição é coerente com dificuldades típicas de estudantes em início de formação.

\subsection{Critérios de avaliação}

Os seguintes critérios foram adotados para avaliação do agente:

\begin{itemize}
  \item \textbf{Consistência quantitativa:} correspondência entre contagem de violações por arquivo no relatório e no gabarito.
  \item \textbf{Cobertura de categorias:} capacidade de representar corretamente os tipos de erro retornados por \texttt{pycodestyle}.
  \item \textbf{Qualidade explicativa:} clareza na descrição do problema e na proposta de correção.
  \item \textbf{Confiabilidade textual:} ausência de afirmações que contradigam os dados de entrada.
\end{itemize}

\subsection{Resultados esperados e próximas etapas}

Espera-se que o agente mantenha precisão total na reprodução das violações reportadas por \texttt{pycodestyle} e entregue feedback didático útil para aprendizado. Também se espera redução de inconsistências textuais no resumo final por meio de prompts com restrições explícitas.

As próximas etapas incluem: (i) adicionar avaliação automática de aderência entre gabarito e saída textual; (ii) ampliar o dataset com códigos reais de projetos acadêmicos; (iii) comparar desempenho da solução com outras ferramentas de lint e formatação; e (iv) investigar integração com interfaces web para uso em monitoria.

\section{Resultados, Ética e Discussão}

\subsection{Resultados obtidos}

Na execução do pipeline sobre os 25 arquivos do dataset, o \texttt{pycodestyle} identificou 217 violações distribuídas em 12 arquivos. Os 13 arquivos restantes não apresentaram violações. As categorias detectadas foram: E225 (62), E231 (78), E302 (30), E301 (21), E201 (10), E202 (10), E228 (3), E501 (2) e E227 (1).

Com base nesses dados, o agente gerou relatórios por arquivo contendo explicação do erro e sugestões de correção. A inspeção dos resultados indica aderência completa da contagem por arquivo ao gabarito utilizado, evidenciando que a etapa determinística conduz corretamente a análise técnica.

Além da aderência quantitativa, observou-se ganho qualitativo no uso da LLM: as saídas apresentaram linguagem instrucional, exemplos corrigidos e recomendações adicionais para manutenção e legibilidade. Quando sugestões extrapolam o escopo estrito de \texttt{pycodestyle}, elas são tratadas como recomendações complementares e não como erro objetivo.

\subsection{Aspectos éticos e responsabilidade no uso da IA}

Do ponto de vista ético, a solução prioriza minimização de riscos de privacidade ao executar o modelo localmente, evitando envio automático de código-fonte para provedores externos. Essa decisão reduz exposição de propriedade intelectual e aproxima o projeto das diretrizes da LGPD em contextos educacionais e corporativos.

Outro ponto relevante é a responsabilidade sobre o feedback gerado: embora o agente automatize a revisão de estilo, sua saída não substitui revisão humana em decisões arquiteturais, segurança ou corretude funcional. Assim, o uso recomendado é de apoio pedagógico e pré-revisão técnica.

\subsection{Endereços do projeto e vídeo}

Repositório GitHub: \href{https://github.com/vinicius-cappatti/avaliador-boas-praticas}{https://github.com/vinicius-cappatti/avaliador-boas-praticas}

Vídeo no YouTube: link em produção para a entrega final da etapa 2.

\section{Conclusão}

O projeto demonstrou que a combinação entre análise estática com \texttt{pycodestyle} e geração de explicações por LLM é uma estratégia viável para revisão automática de código Python orientada ao aprendizado. O componente determinístico garante rastreabilidade das violações, enquanto o componente gerativo amplia a utilidade pedagógica do resultado.

Os resultados obtidos no dataset de avaliação indicam que o objetivo principal foi alcançado: construir um agente capaz de revisar código Python com base em boas práticas, mantendo consistência com as regras da PEP 8 e fornecendo feedback compreensível. Como contribuição, a solução oferece um pipeline reprodutível, extensível e alinhado a requisitos de privacidade, com potencial de aplicação em disciplinas introdutórias de programação.

\section{Referências bibliográficas}

PLAAT, Aske et al. Agentic large language models: a survey. \textit{Journal of Artificial Intelligence Research}, [S. l.], v. 84, 2025. DOI: \url{https://doi.org/10.1613/jair.1.18675}.

PYCODESTYLE DEVELOPERS. pycodestyle documentation. [S. l.], [s. d.]. Disponível em: \url{https://pycodestyle.pycqa.org/}. Acesso em: 27 maio 2026.

PYTHON SOFTWARE FOUNDATION. PEP 8: style guide for Python code. [S. l.], 2001. Disponível em: \url{https://peps.python.org/pep-0008/}. Acesso em: 27 maio 2026.

SIMÕES, Igor Regis da Silva; VENSON, Elaine. Evaluating source code quality with large language models: a comparative study. In: \textit{SIMPÓSIO BRASILEIRO DE QUALIDADE DE SOFTWARE (SBQS)}, 2024. DOI: \url{https://doi.org/10.1145/3701625.3701650}.

WANG, Lei et al. A survey on large language model based autonomous agents. \textit{Frontiers of Computer Science}, [S. l.], 2024. DOI: \url{https://doi.org/10.1007/s11704-024-40231-1}.

ZHANG, Chaoning et al. A complete survey on generative AI (AIGC): is ChatGPT from GPT-4 to GPT-5 all you need? \textit{arXiv}, 2023. arXiv:2303.11717. Disponível em: \url{https://arxiv.org/abs/2303.11717}. Acesso em: 27 maio 2026.

\end{document}
